\chapter{Le articolazioni dell'arto superiore}

% ---------------------------------------------------------------------------
\section{Generalità e cinto scapolare}

Lo \textbf{scheletro appendicolare} comprende le ossa delle estremità
superiore ed inferiore, gli arti, e gli elementi di supporto che collegano
gli arti al tronco, detti cinti o cingoli: il \textbf{cinto scapolare} (o
toracico), a cui si articolano gli arti superiori, e il \textbf{cinto
pelvico}, a cui si articolano gli arti inferiori.

Lo scheletro dell'arto superiore comprende il cinto scapolare, formato
dalla \textbf{clavicola} e dalla \textbf{scapola}, e l'arto libero
superiore, a sua volta suddiviso in braccio (\textbf{omero}), avambraccio
(\textbf{ulna} e \textbf{radio}) e mano (polso, palmo e dita, formati
rispettivamente da carpo, metacarpo e falangi).

Le articolazioni dell'arto superiore sono: l'articolazione
sterno-clavicolare, l'articolazione acromio-clavicolare, l'articolazione
glenomerale, l'articolazione del gomito, l'articolazione radio-ulnare
distale e le articolazioni della mano.

Il cinto scapolare è formato da quattro ossa: due clavicole, poste
anteriormente, e due scapole, situate posteriormente; il cinto scapolare
non si articola con la colonna vertebrale. Le articolazioni che
coinvolgono clavicola e scapola sono l'articolazione sternoclavicolare,
l'articolazione acromioclavicolare e l'articolazione scapolotoracica (uno
spazio di scorrimento funzionale tra scapola e gabbia toracica); a queste
si aggiunge, tra scapola e omero, l'articolazione scapolo-omerale (o
gleno-omerale), delimitata superiormente dallo spazio sottoacromiale.

% ---------------------------------------------------------------------------
\section{Articolazione sterno-clavicolare}

L'articolazione sterno-clavicolare unisce l'estremità sternale della
clavicola all'incisura clavicolare del manubrio dello sterno. Lo sterno è
costituito dal manubrio, in alto, dal corpo, al centro, e dal processo
xifoideo, in basso; il manubrio presenta superiormente l'incisura
giugulare, tra le due incisure clavicolari, mentre l'angolo sternale segna
il passaggio tra manubrio e corpo. Nella stessa regione si trovano anche
l'articolazione acromioclavicolare, l'acromion, il processo coracoideo e
la cavità glenoidea della scapola, in rapporto con la prima costa.

L'articolazione sterno-clavicolare è tenuta insieme da tre legamenti
principali. Il \textbf{legamento inter-clavicolare} collega le estremità
sternali delle due clavicole, attaccandosi all'incisura giugulare del
manubrio sternale. Il \textbf{legamento costo-clavicolare} ha la forma di
un cono rovesciato ed è formato da due fasci, uno anteriore e uno
posteriore: il fascio anteriore si dirige dall'alto in basso e
medialmente, quello posteriore dall'alto in basso e lateralmente. Il
\textbf{legamento sterno-clavicolare} è formato da due bande fibrose: una
anteriore, che si estende dalla superficie antero-superiore dell'estremità
sternale della clavicola all'estremità antero-supero-laterale del manubrio
sternale, fino a raggiungere la prima cartilagine costale, e una
posteriore, meno robusta, che si estende dalla superficie posteriore
dell'estremità sternale della clavicola alla porzione superiore della
faccia posteriore del manubrio sternale. Tra le superfici articolari è
inoltre interposto un disco articolare, mentre l'articolazione stessa
prende rapporto con la cartilagine costale della prima costa e con
l'articolazione sternocostale.

Il movimento consentito dall'articolazione sterno-clavicolare comprende la
proiezione anteriore e posteriore della clavicola (protrazione e
retrazione) e il suo innalzamento e abbassamento (elevazione e
depressione); tali movimenti sono accompagnati, a livello della superficie
articolare, da combinazioni di rotolamento e scivolamento (\textit{roll} e
\textit{slide}) tra clavicola, sterno e prima costa.

% ---------------------------------------------------------------------------
\section{Articolazione acromion-clavicolare}

L'articolazione acromion-clavicolare unisce l'estremità acromiale della
clavicola all'acromion della scapola ed è soggetta a traumi e
lussazioni. I suoi mezzi di unione principali sono il \textbf{legamento
acromioclavicolare} e il \textbf{legamento coracoclavicolare}, quest'ultimo
suddiviso in un \textbf{legamento trapezoide} e in un \textbf{legamento
conoide}; a questi si aggiunge il \textbf{legamento coracoacromiale}, che
insieme all'acromion, al processo coracoideo e al legamento
coracoclavicolare forma l'\textbf{arco coracoacromiale}, posto sopra la
testa dell'omero. Sull'omero si riconoscono in questa regione la grande e
la piccola tuberosità, separate dal solco bicipitale, mentre sulla scapola
si osservano l'angolo superiore, il margine mediale, la faccia costale, il
legamento trasverso superiore della scapola e l'incisura della scapola.

I movimenti permessi dall'articolazione acromion-clavicolare sono
piccoli movimenti di scivolamento, in grado di modificare l'ampiezza
dell'angolo che si forma tra la scapola e la clavicola, e movimenti di
rotazione della scapola in senso orario e antiorario. La scapola può così
compiere, rispetto alla clavicola, una rotazione interna ed esterna e
un'inclinazione (\textit{tilting}) anteriore e posteriore, oltre a una
rotazione verso l'alto e verso il basso.

\subsection{Lesioni dell'articolazione acromioclavicolare}

In caso di caduta sulla spalla o sul braccio esteso si determinano spesso
la lussazione dell'articolazione acromioclavicolare e la lesione dei
legamenti coracoclavicolare e acromioclavicolare. Si distinguono tre
gradi di gravità crescente: lo stiramento dei legamenti, la rottura del
legamento acromioclavicolare e, nella forma più grave, la lussazione
totale dell'articolazione acromioclavicolare, con rottura sia del
legamento coracoclavicolare sia del legamento acromioclavicolare.

Nella lussazione acromion-clavicolare si osserva un rilievo netto della
clavicola, che determina un sollevamento della pelle a livello del piano
della spalla; la pressione con un dito permette di abbassare la clavicola
e di ridurre la lussazione, ma, dal momento che si rilascia la pressione,
la deformità si ripresenta (segno del tasto di pianoforte). Le cause più
frequenti sono il trauma diretto sulla spalla, la caduta sul moncone della
spalla e i traumi sportivi. I segni caratteristici sono il rilievo della
clavicola e il dolore localizzato.

Il trattamento prevede, in prima battuta, mobilizzazioni passive e
attive associate alla terapia manuale per la riduzione del dolore e il
recupero del movimento; segue l'esercizio terapeutico per il recupero
dell'arco di movimento, con attenzione, nelle prime fasi, a lavorare al di
sotto dei 90° di elevazione, e per il rinforzo muscolare; infine, la
terapia strumentale può essere impiegata per l'eventuale riduzione del
dolore e del processo infiammatorio.

\subsection{Legamenti propri della scapola}

I legamenti propri della scapola sono quelli che hanno origine e
inserzione solo sulla scapola stessa: tra questi rientrano il legamento
coracoacromiale, il legamento coracoclavicolare e il legamento trasverso
superiore della scapola. A livello della cavità glenoidea si trovano
inoltre il labbro glenoideo, che ne approfondisce il profilo, il tendine
del muscolo bicipite, la capsula articolare e i legamenti glenomerali
superiore, medio e inferiore, che si inseriscono, insieme al tubercolo
infraglenoideo, sul contorno della cavità stessa.

% ---------------------------------------------------------------------------
\section{Articolazione glenomerale}

L'articolazione glenomerale (o scapolo-omerale) si forma tra i capi
articolari della scapola e dell'omero, rispettivamente la cavità glenoidea
e la testa dell'omero. Per la forma emisferica dei capi articolari in
gioco si tratta di un'\textbf{enartrosi} (articolazione a sfera), dotata di
tre assi di movimento perpendicolari fra loro e quindi di sei movimenti
principali; è l'articolazione con il più ampio range di movimento di tutto
il corpo.

La capsula articolare e i legamenti glenomerali decorrono dall'omero alla
cavità glenoidea e aiutano, in modo determinante, la capsula articolare a
tenere insieme le superfici articolari; tra questi si distingue in
particolare il legamento coracoomerale. Nella regione si riconoscono
inoltre la guaina sinoviale bicipitale, il solco bicipitale, il recesso
ascellare, il collo della scapola e, posteriormente, la fossa
infraspinata e la spina della scapola.

La \textbf{membrana sinoviale} è una sottile membrana di tessuto
connettivo, presente nelle articolazioni, che riveste internamente la
capsula articolare e la parte articolare dell'osso: è a contatto diretto
con la membrana fibrosa sulla superficie esterna e con il liquido
sinoviale sulla superficie interna. Nella regione della spalla si
associano alla membrana sinoviale la borsa sottocoracoidea, il legamento
trasverso dell'omero, il tendine del capo lungo del muscolo bicipite
brachiale, che decorre nel solco bicipitale avvolto dalla guaina sinoviale
bicipitale, e la borsa sottotendinea del muscolo sottoscapolare.

\subsection{Cinematica dell'articolazione gleno-omerale}

L'articolazione gleno-omerale possiede quattro gradi di libertà. Sul piano
sagittale si compiono la flessione, da 0° a 180°, e l'estensione, da 0° a
45°. Sul piano frontale si compiono l'abduzione, da 0° a 120°, e
l'adduzione, da 0° a 30°/45°. Sul piano trasversale si compiono la
flessione orizzontale, da 0° a 140°, e l'estensione orizzontale, da 0° a
30°. Sul piano orizzontale si compiono infine la rotazione interna, da 0°
a 70°, e la rotazione esterna, da 0° a 80°.

% ---------------------------------------------------------------------------
\section{Le articolazioni del gomito}

Le articolazioni del gomito sono tre: l'articolazione omero-ulnare,
l'articolazione omero-radiale e l'articolazione radio-ulnare prossimale.
Nel loro insieme sono formate dall'estremità distale dell'omero e dalle
estremità prossimali di radio ed ulna, e funzionano complessivamente come
un cardine (ginglimo): la maggior parte della stabilità è fornita
dall'incastro tra la troclea omerale e l'incisura trocleare dell'ulna. Le
principali strutture ossee coinvolte sono, sull'omero, l'epicondilo e
l'epitroclea (epicondilo mediale), il condilo omerale, la troclea omerale,
le fosse radiale, coronoidea e olecranica, le creste sopracondiloidee
mediale e laterale e il solco condilotrocleare; sull'ulna, il processo
coronoideo, l'incisura trocleare e l'olecrano; sul radio, il capitello
(testa radiale), il collo del radio e la tuberosità radiale.

\subsection{Articolazione omero-ulnare}

Nell'articolazione omero-ulnare la troclea dell'omero si proietta
nell'incisura trocleare dell'ulna; l'articolazione permette i movimenti di
flesso-estensione dell'avambraccio sul braccio ed è classificata come
\textbf{ginglimo angolare} (articolazione a troclea), con un unico asse di
movimento e quindi due movimenti principali.

\subsection{Articolazione omero-radiale}

Nell'articolazione omero-radiale il condilo dell'omero si articola con la
testa radiale (capitello); si tratta di una \textbf{condilartrosi}.

\subsection{Articolazione radio-ulnare prossimale}

Nell'articolazione radio-ulnare prossimale il capitello del radio si
articola con l'incisura radiale dell'ulna; l'articolazione è classificata
come \textbf{ginglimo laterale} (articolazione trocoide, con un asse di
movimento e due movimenti principali) e permette i movimenti di
prono-supinazione della mano.

\subsection{Stabilità, mezzi di unione e traumatologia del gomito}

L'articolazione del gomito è molto stabile: le superfici ossee si
rapportano in modo da prevenire movimenti laterali e rotatori, e la
capsula articolare è molto spessa e rinforzata da legamenti resistenti.
La superficie mediale è stabilizzata dal \textbf{legamento collaterale
ulnare}, costituito da fibre che dall'epicondilo mediale dell'omero si
inseriscono sul margine mediale dell'incisura trocleare dell'ulna, dal
processo coronoideo fino all'olecrano. La superficie laterale è
stabilizzata dal \textbf{legamento collaterale radiale}, costituito da
fibre che dall'epicondilo laterale dell'omero si inseriscono mediante tre
fasci: il fascio anteriore si inserisce davanti all'incisura radiale, il
fascio medio dietro l'incisura radiale e il fascio posteriore
sull'olecrano. Tra la testa del radio e l'ulna si estende inoltre il
\textbf{legamento anulare del radio}, un fascio fibroso che circonda la
circonferenza articolare del radio e si fissa all'incisura radiale
dell'ulna; a questo si associa il \textbf{legamento quadrato}, che connette
il collo del radio all'incisura radiale dell'ulna.

Dal punto di vista traumatologico, le fratture del gomito rappresentano il
7\% di tutte le fratture: sul totale dei casi, il 33\% delle fratture
interessa l'estremo distale dell'omero, il 10\% l'olecrano e il 35\% il
capitello radiale. I traumi sono solitamente indiretti e vengono generati
per caduta con il braccio in estensione o in flessione.

I movimenti permessi dall'articolazione del gomito sono la flessione e
l'estensione dell'avambraccio sul braccio (nella posizione di flessione
sono consentiti anche modesti movimenti di lateralità) e la pronazione e
la supinazione, rese possibili in particolare dall'articolazione
radio-ulnare prossimale. La \textbf{membrana interossea} è una lamina
fibrosa posizionata tra i margini interossei di radio ed ulna, rinforzata
prossimalmente dalla corda obliqua: separa i muscoli anteriori da quelli
posteriori dell'avambraccio.

\subsection{Nota clinica: sublussazione del capitello radiale}

Nei bambini piccoli è comune una lesione dolorosa causata da una trazione
sull'avambraccio, che si manifesta di solito come un rifiuto nel
movimentare il gomito: si verifica quando il braccio viene tirato
bruscamente verso l'alto con l'avambraccio ruotato verso l'interno,
strappando il legamento anulare, che è attaccato in modo lasso al collo
del radio. Quando il capitello radiale fuoriesce dalla cavità, il
legamento può incastrarsi tra il capitello radiale e il condilo omerale.
In genere la rotazione dell'avambraccio verso l'esterno e la flessione del
gomito riportano il capitello radiale nella sua posizione normale.

% ---------------------------------------------------------------------------
\section{Articolazione radio-ulnare distale}

L'articolazione radio-ulnare distale si forma tra il capitello dell'ulna e
l'incisura ulnare del radio, in rapporto con la superficie articolare
carpale, il processo stiloideo del radio, il processo stiloideo dell'ulna
e il tubercolo dorsale; è classificata come articolazione trocoide
(ginglimo laterale) e partecipa, insieme all'articolazione radio-ulnare
prossimale, ai movimenti di pronazione e di supinazione.

% ---------------------------------------------------------------------------
\section{Le articolazioni della mano}

Le articolazioni della mano sono l'articolazione radio-carpica, le
articolazioni inter-carpiche, le articolazioni carpo-meta-carpiche, le
articolazioni inter-meta-carpiche, le articolazioni meta-carpo-falangee e
le articolazioni inter-falangee. Lo scheletro della mano consta delle ossa
carpali, che formano il polso (tra cui scafoide, semilunare, piramidale e
pisiforme), delle ossa metacarpali, che formano il palmo, e delle falangi,
che formano le dita.

L'articolazione radio-carpica si forma tra la superficie articolare
distale del radio, con il relativo disco articolare, e la fila prossimale
delle ossa carpali; è una \textbf{condiloartrosi} (articolazione
ellissoidale), stabilizzata dalla capsula articolare e dal legamento
collaterale ulnare del carpo.

\subsection{Articolazioni inter-carpiche}

Tra le ossa della fila prossimale del carpo si formano delle
\textbf{artrodie}, tra le facce contigue dello scafoide, del semilunare e
del piramidale; il pisiforme, invece, si articola con la faccia palmare
del piramidale. Tra le ossa della fila distale si formano ugualmente delle
\textbf{artrodie}, tra le facce contrapposte del trapezio, del trapezoide,
del capitato e dell'uncinato. L'\textbf{articolazione mediocarpica}, che si
forma tra le ossa delle file prossimale e distale, ad esclusione del
pisiforme, combina caratteristiche di artrodia e di condiloidea.

I movimenti permessi dall'articolazione della mano sono la flessione,
l'estensione, l'abduzione e l'adduzione della mano rispetto
all'avambraccio (rispettivamente deviazione radiale e deviazione ulnare) e
la circumduzione.

\subsection{Articolazioni carpo-metacarpiche e delle dita}

Le \textbf{articolazioni carpo-metacarpiche} sono artrodie, ad eccezione
dell'articolazione del pollice, che è un'articolazione a sella; permettono
flessione, estensione e una modesta lateralità. Il primo metacarpale è più
mobile degli altri e può compiere movimenti di flessione, estensione,
abduzione e adduzione. Le \textbf{articolazioni inter-metacarpiche} sono
artrodie che si formano tra le basi degli ultimi quattro metacarpali e
consentono soltanto movimenti di scivolamento. Le \textbf{articolazioni
metacarpo-falangee} sono articolazioni a condilo, tra le teste delle ossa
metacarpali e le basi delle falangi prossimali, e permettono flessione ed
estensione e, in misura minore, abduzione, adduzione e circumduzione. Le
\textbf{articolazioni inter-falangee} sono a troclea, tra le teste delle
falangi e le basi delle falangi successive: sono due per ciascun dito, ad
eccezione del pollice, che ne possiede una sola, e permettono ampi
movimenti di flessione ed estensione.
